\subsection{Future Directions}

Four research priorities emerge from this review. First, forecasting needs standardized clinical benchmarks. Detection has the ILAE Phase 3 benchmark of FPR below 0.05 to 0.1 per hour. Forecasting lacks an equivalent standard. AUC values around 0.74 to 0.75 appear consistently across forecasting studies, but the clinical meaning of this performance level remains unclear. Metrics such as IoC and TiW need standardization. The field should establish minimum acceptable performance levels for forecasting devices to enable clinical translation.

Second, large-scale home validation is needed. Only two detection studies include prospective home validation. Hospital and EMU settings may not capture the diversity of real-world activities that cause false alarms. Exercise, cooking, household chores, and other daily activities present challenges not encountered in controlled environments. The 0.165/h FPR reported by \textcite{dongTwoLayerEnsembleMethod2022} in home settings exceeds the clinical benchmark. Future studies should prioritize prospective validation in intended use environments with sufficient sample sizes and duration.

Third, novel modalities are needed for non-convulsive seizures. Current approaches detect only motor seizures with convulsive manifestations. Focal aware seizures without motor signs remain undetectable with current wearable approaches. \textcite{spahrDeepLearningbasedDetection2025} achieved excellent performance for convulsive seizures using accelerometer data alone, but this approach cannot detect seizures without movement. Multi-modal autonomic approaches show promise but have not addressed non-convulsive seizure types. Novel biosignals or advanced feature engineering may be required to detect the subtle physiological changes associated with focal seizures.

Fourth, adaptive personalization strategies could address the trade-off between deployability and individual performance. Global models enable immediate deployment but may not work equally well for all patients. \textcite{meiselMachineLearningWristband2020} achieved 62\% patient success with LOSO validation, indicating substantial inter-patient variability. Patient-specific models achieve higher individual success but require weeks to months of individual data. \textcite{stirlingForecastingSeizureLikelihood2021} required a mean of 14.6 months per patient. Mixed approaches that combine population-level patterns with individual adaptation, such as the method used by \textcite{vielufSeizureMonitoringCombined2025}, may offer a balance. Research should address the cold-start problem by developing models that can initialize from population data and adapt efficiently to individual patients.
